\section{第 5 章 关键功能模块设计与实现}

在对各个核心业务模块进行展开之前，首先要说明首页导航页在系统中起到的作用。不同于单纯的静态欢迎页，本系统首页具有“统一入口、实时概览、模块分发”这三种功能。一方面，首页用统一的导航栏和功能卡片把多模式预报、探空分析、历史气候、农业气象等模块组织到同一个视觉框架里，使整个平台的结构更清晰；另一方面，首页顶部实时天气区和智能提醒区又能提前给用户在进入专业分析页面之前提供当前天气状态、体感信息以及重点提醒，从而提高平台整体的连贯性和可用性。因此首页并不是全文的独立技术主线，但是它是系统页面组织和业务引导中不可缺少的一环。

\subsection{5.1 多模式预报模块设计与实现}

\begin{figure}[htbp]
\centering
\caption{图 5-1  多模式预报模块流程图}
\label{fig:5-1}
\end{figure}

\begin{figure}[htbp]
\centering
\caption{图 5-2-1  多模式预报页面示意图}
\label{fig:5-2}
\end{figure}

\begin{figure}[htbp]
\centering
\caption{图 5-2-2  多模式预报页面示意图}
\label{fig:5-2}
\end{figure}

多模式预报模块属于本文系统的重要模块之一，也是本文业务主线中最为重要的部分。不同于传统的天气页面只给出一个预报结果，本模块更加强调两个问题，一是怎样使未来天气的变化更加连续、直观，二是怎样让用户看到不同的模式之间的差别，从而理解预报结果的不确定性以及参考价值。

按照此思路，系统并没有把预报功能设计成单一的列表，而是分成了两个互相配合的层次。第一部分为72小时精细化预报，主要给出未来三天温度、降水、湿度、风速在各时间段的变化情况，适合普通用户做短期出行和生活安排使用。第二部分为多模式对比分析，主要用以展示 ECMWF、GFS、ICON 等模式在同一天线上的差异，从而让使用者看到各个模式对于天气过程强弱、转折时间及趋势发展是否一致。用这样的双层结构来保留日常使用方便的同时又提高页面分析能力。

技术实现上，本模块使用Flask作为页面和接口的载体，用Open-Meteo API作为主要的预报数据来源，在服务层对多模式结果进行统一组织。前端使用Plotly.js画图，时间序列温度、降水、风速等可以用交互式的曲线来表现。与静态表格相比，此种方式更能体现连续的天气过程，也更有利于在同一个坐标系中比较不同模式之间的差异。

为了加快页面响应速度，在该模块中使用了短时缓存以及异步加载的方式。页面在第一次加载的时候先完成整体结构的渲染，然后分别请求精细化预报和多模式数据，从而避免由于一个接口响应较慢而阻塞整个页面。同时短周期缓存可以降低对于同一个预报数据重复请求的情况，在保证结果时效性的基础上提高用户的体验。对多模式场景系统采用并行获取不同模式结果的方式来减小等待时间，使页面更适合作为在线访问。

本模块除了原始预报的显示之外，还有对机器学习误差校正的主要应用环境。系统把 ECMWF 预报链路同误差校正模型联系起来，用户在查看逐小时预报的时候可以决定是否采用校正结果。该种设计不是把机器学习模块作为一个独立的页面来实现，而是在预报功能的基础上，将机器学习模块作为一个增强功能嵌入到主业务流程中。因此用户既可以用原始预报为基准，也可以在需要的时候查看订正后的温度、湿度、风速和降水结果，使用体验更自然。

需要指出的是，系统并没有把机器学习校正无差别地应用到所有的模式上，只在 ECMWF 这一条主链路上启用。该处理符合训练口径和运行口径一致的原则，也可以防止把同样的校正结果机械地套用到不同的模式上所造成的误导。同时系统还给校正结果设计了回退机制，当模型不可用或者运行时的条件不符合要求的时候，页面仍然可以稳定地返回原来的预报结果，保证主功能不受影响。

总体上来说，多模式预报模块把 Flask 页面组织、Open-Meteo 数据调用、Plotly.js 图表展示、短时缓存、异步加载、机器学习误差校正有机地结合在一起，不但提高了天气信息的可读性，而且使系统由单纯的预报展示平台发展成为具有一定分析性质的综合气象服务平台。

\subsection{5.2 探空分析模块设计与实现}

\begin{figure}[htbp]
\centering
\caption{图 5-3  探空分析模块流程图}
\label{fig:5-3}
\end{figure}

\begin{figure}[htbp]
\centering
\caption{图 5-4  探空分析页面示意图}
\label{fig:5-4}
\end{figure}

探空分析模块是本系统中专业性最强的功能模块之一。相比地面天气预报主要回答“未来天气怎么样”，探空分析更关注“大气当前处于什么状态、是否具备不稳定条件、是否存在对流发展潜势”。因此，这一模块的作用并不只是补充一个新的数据页面，而是承担了平台专业分析能力展示的重要任务。多模式预报模块解决的是“未来会发生什么”，而探空分析模块解决的是“当前大气为什么会这样，以及是否具备进一步发展的环境条件”。两者结合后，系统不仅能提供面向公众的天气结果，也能提供一定程度的气象诊断依据。

从业务需求出发，系统引入探空模块主要是为了解决三个问题。第一，单纯查看地面温度、降水和风速，并不能完整判断对流潜势、逆温层结构或高空风场特征，特别是在雷雨、短时强降水和强对流天气背景下，仅靠地面信息往往不足以支撑分析。第二，原始探空资料通常以文本和层结表形式提供，专业性较强，普通用户难以直接阅读，因此需要将其转换为图形和关键参数。第三，系统不仅要为专业用户保留探空图和指数，也要为普通用户提供风险等级、通俗解读以及面向实际场景的影响提示。

基于这一思路，探空分析模块被设计成一个集“探空资料获取、垂直结构分析、专业图形展示、关键指数提取和结果解释”于一体的综合诊断模块，而不是单纯的数据下载功能。当前系统的探空数据主要来自怀俄明大学在线探空资料服务，后端通过 requests 请求原始资料，并使用 Pandas 对层结信息进行整理，再交给分析和绘图模块处理。图形部分主要依赖 Matplotlib 和 MetPy 完成专业气象图生成，页面端则采用 Bootstrap 布局和前端 JavaScript 异步更新方式完成展示。

从页面结构上讲，探空分析模块是用左侧控制、右侧展示的方式来安排的。左侧为时间选择、图形类型切换、关键指数显示，右侧为探空图、风险评估、分析说明。该结构使用户可以按照“先选时次、再看图形、再看结论”的顺序来完成阅读，更符合探空资料使用的习惯。目前模块主要有四种图形，T-lnP图、Skew-T图、高空风矢图、卡通直观图。T-lnP图、Skew-T图主要用来表示温度、露点、风场、大气稳定结构；高空风矢图反映不同高度的风速、风向的变化；卡通直观图把专业结果更易理解地转为可视化的形式，方便非专业用户快速掌握大气状况。

除了图形以外，系统还会同步显示一组重要的指数，即CAPE、CIN、K指数、抬升指数、0-6km垂直风切变和整层可降水量。这些指标被集中到参数区，配有简要说明，方便用户在没有完全依靠探空图的情况下先做初步判断。CAPE反映对流有效位能，CIN反映抑制能量，风切变与风暴组织化程度有关。对于非专业用户来说，图形和参数的结合使用比单独看探空图更容易理解；对有基础的用户来说，这些指标可以很快找到分析的重点。

探空模块的另一个特点就是除了给出参数、图形外，还会给出风险等级以及场景化的解释。系统目前会将探空结果进行归纳，给出较为通俗的解释，尝试从航空、农业、日常生活等角度来补充影响提示。该种设计使探空模块不再只是专业图件的展示页面，而是可以将探空分析的结果转化为实际应用信息的综合模块。对航空来说，探空资料可以用来判断层结稳定性、风切变以及对流风险；对农业来说，可以理解强对流、短时强降水或者逆温状况给农业带来的影响；对公众用户来说，可以转化为比较直观的天气风险提示。

探空模块最值得强调的支撑设计之一，就是对时间口径进行统一处理。由于页面面向国内用户，时间输入以北京时间理解更自然；但是探空业务资料一般用00Z、12Z标准UTC时次组织，所以系统在后端会先把北京时间转换成UTC，再规整到最近的标准探空时次。若目标时次没有有效资料，系统可以回退到前一个标准时次，将最终实际分析所用的时次反馈给用户。该机制本身不是模块的主要卖点，但是它可以保证探空分析在实际使用过程中是稳定的、可解释的。

另外，系统对于参数处理更看重结果的可靠性而不是形式的完整。比如在CAPE进行补充计算的时候出现负值的情况，页面不会给出一个可能会误导用户的数值，而是一直用N/A来表示异常。与把所有的指标都用数字来表示的思想不同的是，这样一种处理方式更符合面向真实用户的应用系统的要求，并且有助于减少错误结论的流传。

探空分析模块把怀俄明大学探空资料、requests数据抓取、Pandas结构化处理、Matplotlib和MetPy专业绘图、Bootstrap+JavaScript页面交互结合起来，使系统在多模式预报的基础上又有了垂直大气结构诊断和场景化解读的能力。该模块明显加强了整个平台的专业程度，也是论文中能够体现系统“分析能力”而不是“展示能力”的部分之一。

\subsection{5.3 历史气候分析模块设计与实现}

\begin{figure}[htbp]
\centering
\caption{图 5-5  历史气候代表事件识别流程图}
\label{fig:5-5}
\end{figure}

\begin{figure}[htbp]
\centering
\caption{图 5-6  历史气候年度页示意图}
\label{fig:5-6}
\end{figure}

历史气候分析模块是本系统中最具统计分析特征的功能模块之一。与多模式预报关注未来几天、探空分析关注当前时次的大气结构不同，历史气候模块更强调从多年观测样本中提取背景信息，帮助用户理解某一年的整体气候特征、长期变化趋势以及典型异常事件在历史中的相对位置。正因为有了这一模块，系统才能把“短期天气”“当前环境”和“长期背景”连接起来，形成更完整的分析链条。

从业务需求出发，用户在看到某一年的高温、暴雨或低温事件时，通常不仅关心它本身的数值大小，还会进一步关心这一年整体上是偏暖还是偏冷、偏湿还是偏干，以及某个异常事件到底只是绝对值较大，还是在历史背景中也具有代表性。基于这一需求，系统将历史气候分析组织为三条主线：年度气候概览、长期趋势分析和代表事件识别。这样的设计使历史气候模块不再只是“看历史数据”，而是能够从年度统计、多年趋势和典型事件三个角度理解气候特征。

在页面结构上，历史气候模块当前分为两个主要页面，即年度分析页 /history/yearly 和趋势分析页 /history/trend。年度分析页主要反映某一年的平均气温、总降水量、雨雪日数、高温日等指标，与标准气候态进行比较；趋势分析页将多年年度统计结果连贯起来，用来观察气温、降水、极端天气频次在更长时间尺度上的变化规律。双页面结构既可以满足用户对某一一年的细看需求，也可以满足用户对多年总体变化的把握。

从技术实现角度来讲，历史气候模块用 Flask 作为页面入口和接口组织形式，分析层主要是使用 Pandas 和 NumPy 来完成历史观测数据的解析、清洗、按日聚合和统计计算。可视化部分用Plotly、Plotly Express和前端Plotly.js完成年度页和趋势页中交互式图表的展示，比如月度对比图、温度分布图、气温趋势图、降水趋势图和极端天气频次图等。风玫瑰图同样用极坐标图的方式表现出来，使用户可以更加直观地看出年度或者月度风况的分布情况。

历史气候模块的第一个重要功能，是年度气候概览。系统会对指定年份的年平均气温、年降水量、雨日、雪日、高温日、低温日等指标进行统计，并与标准气候态进行对比，帮助用户快速判断该年份是偏暖、偏冷还是偏湿、偏干。当前系统支持两套气候态口径，即 1981-2010 和 1991-2020。允许用户在不同气候态之间切换，可以提高结果的可比性，也使模块更符合气候分析业务的实际表达方式。

第二个功能就是长期趋势分析。系统把多年的年度统计结果串联起来，主要体现气温变化趋势、降水变化特点、高温日、低温日、暴雨日、雷暴日、降雪日等极端天气频次的改变状况。毕业设计论文中这一部分内容尤为重要，它使模块具有了较强的分析性、研究性，而不仅仅是某一年份的静态说明。

第三个也是最具有特色的功能就是事件识别。系统在这一步并没有只保留传统的“原始极值榜单”，而采用的是“原始极值+EW指数”双模式的设计。原始极值模式是以绝对数值的大小来排列的，如年内最高温、最低温和最大日降水量等，直观易懂，所以被当作默认的显示方式。但是只看绝对值的大小不能完全体现一个事件在历史背景中所具有的特殊性，所以系统又引入了EW指数模式。

EW 指数模式的核心思想是：将某一日的高温、低温或降水事件放回历史同月同日样本中比较，观察它在历史背景中所处的位置。对于高温和降水，百分位越高说明越异常；对于低温，则采用反向百分位处理，使越低的温度对应越高的极端程度。为了避免仅凭同日样本就把绝对值不强的事件抬上榜单，系统还加入了同季绝对门槛限制。与此同时，为避免连续高温、连续低温或持续降水过程在榜单中重复刷屏，系统又进一步采用了过程去重思路，将相邻异常日合并为一个过程，并仅保留其中最具代表性的一天。最终，EW 结果按春、夏、秋、冬四季分组展示，而不是全年混排，从而避免某一季节占满全部榜单。

这套“历史同日背景比较+同季门槛+过程去重+季节分组”的设计，把历史气候模块由简单的图表展示提升到代表事件识别的层次上。对用户来说，它能够更好地回答“这一年哪些天气过程最值得被记住”，对于论文来说，它也是一份很好的模块分析亮点。

从总体上讲，历史气候分析模块把历史观测数据统计、气候态对比、长期趋势分析、EW代表事件识别、Plotly系列图表展示有机地结合在一起，具有较强的长期背景分析能力。它既弥补了平台对于时间尺度的不足，又加强了系统对于异常天气以及气候变化问题的解释力度。

\subsection{5.4 机器学习误差校正模块设计与实现}

\begin{figure}[htbp]
\centering
\caption{图 5-7  机器学习误差校正流程图}
\label{fig:5-7}
\end{figure}

机器学习误差校正模块与多模式预报模块的关联最强，它并不是一个孤立存在的算法实验功能，而是直接服务于 ECMWF 预报结果质量提升的增强链路。对用户而言，真正关心的仍然是“未来天气会怎样”，而这一模块所做的事情，是在原始预报结果基础上尽量减少系统性偏差，使页面展示的温度、湿度、风速和降水更贴近真实观测。

从业务背景来看，结合江南地区特别是杭州样本的使用体验和本系统训练、测试样本的表现，可以看到 ECMWF 在中远期时效上往往存在一定的偏暖、偏湿以及降水量级偏差倾向。这种系统性误差并不意味着模式本身不可用，而是说明其原始输出在落到具体站点业务场景时，仍然有必要结合历史样本进行后处理订正。也正因为如此，本模块的重点并不只是“用了机器学习”，而是“如何把模型真正接入预报业务链路，并让结果更可信”。

在数据基础上，当前系统的训练样本主要来自 data/hangzhou\_openmeteo 目录下持续采集的新版 ECMWF 预报文件。按照训练清单目前共对100个预报文件进行了审计，其中97个进入到训练流程当中，另外三个由于采集的时间不符合支持窗口的要求而被剔除出去。训练和运行时都使用杭州国家站/馒头山作为目标站点，坐标为30.2444，120.1528。该种统一的动作十分关键，它很好地解决了旧方案里预报点、训练样本以及观测点的口径不完全相同所造成的难题，使得模型的输出更加具有业务价值。

在样本组织上，系统围绕 issue\_time 进行统一处理。对于新版 HZ\_forecast\_*.csv 文件，凌晨 04:15 左右采集的文件映射到前一日 20:00 本地起报，对应 12Z 周期；下午 16:15 左右采集的文件映射到当日 08:00 本地起报，对应 00Z 周期。训练、验证和测试的划分也以 issue\_time 为主线，而不是对所有目标时刻随机打散，从而避免信息泄漏。与此同时，运行时预报页面也会显式传入 issue\_time，使训练与推理采用的是同一套时次口径。

系统用scikit-learn提供的随机森林模型来建立订正链路。温度、湿度、风速用回归模型做误差订正，降水用两阶段法，先判断有无降水，再估计有降水时的降水量。特征设计的时候系统地保持了预报值本身、时间特征（小时、月份、工作日/周末）、时效特征（lead\_hours、lead\_category、周期性时间编码）这些。由此可以看出本模块并没有一味地追求复杂的模型，而是在目前的业务场景下选择了更合适、也更易于和运行时接口进行对接的方案。

在业务接线上，机器学习误差校正模块并没有单独做成一个独立页面，而是直接嵌入多模式预报模块，作为 ECMWF 预报结果的增强能力出现。当前前端在启用 ML 校正时，会将逐小时 ECMWF 预报数据与运行时估计的 issue\_time 一并提交给 /apply-ml-correction 接口；后端完成校正后，再将温度、湿度、风速和降水的订正结果返回页面。如果模型已成功加载，则页面显示校正结果；如果模型未加载、时次不明确或特征条件不满足，则系统自动回退到原始预报值。这样一来，模型不会打断预报主链路，而是以“可开关、可回退”的方式增强 5.1 模块。

当前系统不仅完成了模型训练，还同步生成了 training\_manifest.json 和 metrics\_report.json 两类追踪产物。前者记录训练数据目录、站点信息、文件审计情况、样本时间范围，后者记录连续变量和降水分类/回归的交叉验证和最终测试结果。这就使得模型训练不再只是单次的离线实验，而成了可以追溯、可以复现、可以不断迭代的一个业务过程。

为了更好的体现两版模型在特征设计方面存在的不同，本文将V1和V2在特征设计上主要的不同点归纳成表5-1。

V2是在V1的基础上，加入了季节、起报周期、短期变化量这些特征以更好的表现误差的时序依赖以及季节性规律。相比于只依靠基本预报值和通用时间变量，这类增强特征更加重视同一模式误差随着季节、起报时次和预报时效一起变化的业务事实，因此更适合用作后面boosting实验方案的输入基础。

表 5-1  V1 与 V2 机器学习误差校正特征差异对比

从最终测试结果看，温度、湿度和风速都表现出不同程度的改善，其中温度和风速的提升更为明显。下表给出连续变量在最终测试集上的主要评估结果。

表 5-2  V1 误差改进幅度

对于降水，系统同时给出分类结果与量级结果。需要诚实说明的是，降水校正并非在所有分类指标上都优于原始预报，但在降水量级误差上改善较为明显，这也说明降水后处理的难度高于温度和风速，不能只依赖单一指标下结论。下表给出降水在最终测试集上的主要评估结果。

表 5-3  V1 降水误差与原始版对比

这一结果是可以在论文中作出合理解释的。降水本身就具有更强的随机性、局地性，即使是在同一个模式、同一个站点口径下，降水是否出现也会比温度、风速更难被准确刻画。其次，目前校正模型策略上更多地减少了“报了雨但实际上没有下雨”的误报，因此Recall有所降低，Precision基本不变或者略有提高，说明模型整体上更保守。从业务应用角度来讲，降水量级误差明显降低有实际意义，当系统判定有降水的时候，校正后的雨量值一般更接近真实的状况，对农业灌溉安排、水库泄水调度、田间排涝准备等场景来说，比起只注重提高降水召回率而言，更有实际的应用价值。

需要指出的是，本模块的意义并不在于保证所有变量、所有折次都一定优于原始预报。事实上，交叉验证中的不同折次仍可能出现改进幅度波动，甚至个别折次变差的情况。但在更统一的站点口径、更统一的时次口径以及更严格的评估条件下，这一版方案显然比旧版更真实，也更具参考价值。对于论文而言，这种“结果更可信”本身就是非常重要的结论。

因此，机器学习误差校正模块把scikit-learn、Pandas、NumPy、Flask、训练清单和评估报告有机地结合在一起，使得系统在多模式预报的基础上又有了后处理优化的能力。它不但提高了主要预报结果的实用性，而且使整个平台工程完整性、技术深度都得到了提高。

\subsection{5.5 农业气象服务模块设计与实现}

\begin{figure}[htbp]
\centering
\caption{图 5-8  农业气象服务模块流程图}
\label{fig:5-8}
\end{figure}

\begin{figure}[htbp]
\centering
\caption{图 5-9  农业气象总览页示意图}
\label{fig:5-9}
\end{figure}

农业气象服务模块是本系统面向行业场景延伸的重要组成部分，其核心目标不是简单罗列若干作物信息，而是把前文已经构建的天气预报、指标计算和风险识别能力，进一步转化为面向农业生产的管理建议。相较于公众天气服务更关注“未来天气如何变化”，农业场景更关心“这种天气会对当前作物阶段产生什么影响，以及接下来该做什么”。因此，本模块围绕“作物状态识别、环境适宜度判断、风险预警与农事建议”四条主线组织业务逻辑，使系统能够从通用气象平台进一步延伸为具有一定决策支持能力的农业气象应用。

在业务对象上系统选择了浙江地区具有代表性的一些作物，即单季晚稻、西湖龙井、杨梅、柑橘和小白菜，并且给每种作物建立了包含生育阶段、适宜温度、水分需求、积温参数以及气象风险阈值的作物知识库。不同作物的管理重点也不一样，茶叶更注重倒春寒、高温灼伤，杨梅更注重成熟期暴雨、大风、高湿引起的落果烂果风险，柑橘更注重冻害、夏季日灼，小白菜则具有周年可种、周期短等特点。将业务知识结构化后，在天气数据出现的时候，系统不仅可以给出气温多少、降水多少等天气情况的描述，还可以进一步解释这些天气条件给不同的作物带来了怎样的影响。

在页面结构上，农业模块由总览页 /agro-dashboard 和作物专题页 /crop/<crop\_id> 组成。总览页是全局农业气象研判的入口，集中显示风险提示滚动条、作物实时状态列表、未来七天农事日历、风险预警清单和全局农情研判结果，使用户在开始使用时就对近期农业生产环境有一个大致的认识。作物专题页进一步细分到单个作物上，呈现当前生育阶段、积温进程、阶段性风险阈值、未来 7 天环境变动以及个性化的农事提议，让用户可以针对某个作物展开更为精确的管理判断。该种“总览+专题”的双层结构，既可以满足管理者快速浏览整个情况的需求，也可以满足种植者对单个作物具体建议的详细查看。

从业务处理流程看，农业模块以上游预报模块输出的未来天气数据为基础，由后端先将逐小时预报转换为逐日农业气象信息，再进入作物分析流程。系统在 build\_agro\_forecast\_list() 中按日聚合温度、降水、湿度和风速，并进一步计算 ET0、净有效降水 precip\_eff、浅层土壤湿度和短波辐射等指标。这样做的意义在于，农业场景并不总是直接关心小时级的原始天气值，而更关注“未来几天总体偏湿还是偏干、适不适合作业、是否存在持续性不利条件”。通过把天气数据转换成更贴近农业语言的日尺度指标，系统后续的风险判断和建议输出也就更容易与实际管理需求对接。

系统根据日期来判断各个作物目前所处的生育阶段，再计算出各个作物的适宜度评分以及风险预警。适宜度评分由温度、水分、极端天气三者共同决定，得到的综合分值再用在总览页上直观体现各种作物目前所处的环境状况。风险预警是由 agro\_alert\_engine.py 来驱动的，它把未来七天每天的预报和作物阈值进行对比，找出低温冻害、高温热害、暴雨渍涝、大风落果、高湿病害等风险，并且也发现了适合喷药、施肥、修剪等田间作业的时间窗口。因此模块输出的不单单是被动的风险提醒，还有“可利用的天气机会”，更加符合农业管理“趋利避害并重”的实际需要。

积温分析属于本模块的重要功能。根据不同的作物设置基点温度和目标积温，在作物专题页中给出当前积温、未来7天的预测新增积温和总体进度百分比，供作物发育进程及成熟节律参考使用。历史积温部分采用ERA5-Land再分析日值数据，计算得到的当前实现是根据指定起始日期到昨日的日最高、日最低气温，然后转换成累计积温；未来 7 天新增积温直接使用系统当前主链路中的 ECMWF 预报结果进行估算。因此选择这样的组合，一方面是因为农业场景下连续、长期、稳定的单站历史逐日观测很难完整获取，再分析资料在时间连续性以及历史覆盖范围上更符合支撑整季积温累计的要求；另一方面，积温属于长期积分量，日尺度误差在跨月甚至跨季累积时会被一定程度平滑，所以ERA5-Land用来做历史积温估算是合理的。相比之下，未来7天中部分更加重视对当前业务预报的响应，所以仍然使用ECMWF作为前向估计的依据，也可以保持和平台预报主链路一致。对小白菜等短周期作物来说，页面可以按照用户的播种日期来计算出合理的估计。积温功能的价值就在于它把天气条件和作物生长过程联系得更加紧密，使用户看到的不只是单次天气的影响，而是与作物发育进度直接相关的连续过程判断。

就建议生成而言，系统采取“AI 建议 + 规则回退”并行的办法来改善可用性。总览页的全域农情研判是由farming\_ai\_adviser.py结合预警信息、未来天气、作物状态生成的，作物专题页是通过/api/agro/advice接口为单个作物请求个性化的农事建议。如果DeepSeek API可用，系统就会给出更加自然、更加完整的农事说明；如果接口不可用、返回异常或者解析失败，那么就直接采用规则引擎给出的结构化建议。该种设计可以防止模块对单个AI服务产生硬依赖，在真实的运行环境中也可以稳定地给出可执行的提示信息。前端页面使用 Flask 模板、Bootstrap 组件、JavaScript 交互和 Chart.js 图表来实现总览展示、AI 刷新、环境曲线可视化等功能，使农业模块具有较好的业务完整性以及页面表达能力。

从应用价值来说，农业气象服务模块不是全文技术主线，但是也是系统场景落地能力最强的部分之一。它表明本系统并不只是简单的天气展示平台，而可以将气象要素同作物生长、田间管理、灾害防范联系起来，产生更加贴近实际生产的服务输出。对于茶园管理、果园病害防治、水稻抽穗扬花期防灾、蔬菜短周期种植安排等场景，以天气分析为基础、以作物管理为落脚点的服务方式有较强的实用性。该模块既可以体现系统应用的拓展性，也可以体现预报分析能力向行业服务转化的完整闭环。

